\chapter{Interaction as a contrast of probabilities}\label{ch:probability-interactions}

A canonical density gives the coalition table a second language: endpoint values become log density ratios, and interactions become alternating log contrasts. The same Möbius operation carries the account across this bridge. We also distinguish these endpoint contrasts from probabilities of choosing actions, so that physical dependence is not confused with a behavioural rule.


\section{The same table in a second language}
Under the canonical density relation, a decrease in potential corresponds to an increase in equilibrium density. Each coalition value can therefore be read as a log density ratio. Applying the same alternating signs to the log densities gives the same interaction coefficient.

This is a continuity result: the object being transformed has changed language, while the finite coalition operation has stayed the same. It is not a new source of value to add to the physical account. When its assumptions hold, the two expressions describe one comparison.

\begin{theorem}{Canonical-density Möbius identity}\label{thm:probability-mobius}
Assume a complete coalition table, one fixed reference measure and support, $0<Z<\infty$, and positive finite densities $p(x_S)=Z^{-1}e^{-V(x_S)}$ at all required endpoints. Then
\begin{align}\label{eq:probability-value}
 E(S)&=\ln\frac{p(x_S)}{p(x_0)},\\
 m_E(T)&=\sum_{S\subseteq T}(-1)^{|T|-|S|}\ln p(x_S)
 \qquad(T\ne\varnothing).\label{eq:probability-mobius}
\end{align}
Equivalently, the coefficient is the logarithm of the product of densities raised to the corresponding alternating powers.
\end{theorem}
\begin{proof}
The density relation gives $\ln p(x)=-V(x)-\ln Z$. Subtracting the two endpoint logs gives Equation~\ref{eq:probability-value}. A nonempty Boolean difference annihilates both constant normalization terms. The remaining signed log sum is the EBU coefficient. Exponentiating the sum gives the product form.
\end{proof}

The algebra uses the density relation at the named points. To call it a physical canonical result, the ensemble and calibration conditions of Chapter~\ref{ch:equilibrium-bridge} must be justified independently. Neither stationarity nor reversibility is proved by this finite algebra.

\section{The pair product}
Writing $p_S=p(x_S)$ and $p_0=p(x_0)$, a normalized two-action table gives
\begin{equation}\label{eq:pair-log-product}
 m_E(AB)=\ln\frac{p_{AB}p_0}{p_Ap_B}.
\end{equation}
The denominator combines the two standalone density readings. The numerator combines the joint endpoint with the shared baseline. The ratio equals one precisely when this pair contrast is zero.

For the courier table, potential values are $16,9,9,4$. Their densities have a common unknown normalization. Substitution gives
$\ln[e^{-4}e^{-16}/(e^{-9}e^{-9})]=-2$, with the normalizer canceling. The normalizer cancels exactly. The same pair correction has been expressed in density language, providing two compatible readings of one declared landscape.

\section{The triple product}
The eight-term contrast becomes
\begin{equation}\label{eq:triple-log-product}
 m_E(ABC)=\ln
 \frac{p_{ABC}p_Ap_Bp_C}{p_{AB}p_{AC}p_{BC}p_0}.
\end{equation}
The same recursive interpretation survives: it measures how a pair log contrast changes when the third action is present. The powers and signs are fixed by Boolean inversion. They must not be guessed from a verbal claim that synergy should be positive.

When densities are very small, direct products can be numerically unstable. A future implementation would normally add log densities with their signed coefficients. That is a numerical representation choice, not a change in the physical definition. Zeros require separate treatment: a zero-density endpoint produces an undefined or infinite logarithm, not a finite coefficient to be smoothed away without a model.

\section{Endpoint densities are not action probabilities}
The entries $p(x_S)$ describe density at physical endpoints. They are not generally probabilities of the binary action pattern $S$. Different coalitions can reach the same endpoint, and a physical density need not sum to one when sampled at a handful of endpoints.

One can define artificial normalized weights $q(S)=p(x_S)/\sum_Rp(x_R)$ on the action patterns. Their nonempty log contrasts match the endpoint contrasts because this new normalizer also cancels. But that construction does not show that actors actually choose according to $q$, or that the resulting binary variables are causally independent.

If a positive binary joint distribution $q$ is independently specified, its pair contrast on a fixed face is an ordinary log odds ratio. A zero pair contrast in the face with all other variables absent does not establish unconditional independence. For example,
\begin{equation}
 \ln q(z)=\text{constant}+a z_1z_2z_3
\end{equation}
has zero pair contrast for actions 1 and 2 when $z_3=0$, but pair contrast $a$ when $z_3=1$. The context changes the relationship exactly as the recursion predicts.

\section{Several quantities called information remain different}
A pointwise log-density contrast is not an alternating sum of marginal entropies. Interaction information uses marginal entropy combinations; multi-information compares a joint distribution with the product of its marginals. Connected information compares successive maximum-entropy models. Cumulants use a log moment-generating function and a different combinatorial structure.

These constructions can all describe dependence, but they answer different questions and need not give the same number or sign. A claim that EBU interaction is one of them requires a precise derivation under the relevant distribution and coding convention. The common word ``information'' does not establish the equality.

For EBU, precise names keep each capability useful. A pointwise contrast explains the endpoints of a particular intervention family; an entropy statistic summarises a distribution. Choosing the quantity that matches the service question avoids a score whose meaning changes as it moves between reports.

\section{The prior-art connection}
Gibbs and log-linear representations, Möbius inversion, cooperative-game dividends and interaction decompositions have extensive antecedents. Jansma explicitly connects higher-order energy interactions, probability contrasts and game-theoretic decompositions in this general mathematical family \cite{jansma2025,jansma2023}. The EBU formula here should be read in that context, not presented as a newly invented inversion principle.

The derivation above fixes its own sign from $E=V(x_0)-V(x_S)$ and $\ln p=-V-\ln Z$. This provides a direct check without relying on an isolated intermediate equation in another text. The versioned equation-level source comparison remains available in the project's prior-art appendix \cite{smg-prior}.

\section{Two expressions for one physical comparison}
The physical-potential expression and the log-density expression agree under the same canonical assumptions. Their equality is a theorem, not an extra value to add to the account. If $E(\varnothing)$ is nonzero because the background evolves, every nonempty contrast still agrees: the raw-to-relative shift is constant across the table.

This gives EBU a consistent statistical representation of shared effects. Pair and triple contributions can be explained through physical endpoint values or through canonical density contrasts, with the same sign convention and reference measure. The representation can reveal a missing normalization factor or an incompatible state description before an institutional rule uses the number.

For actual measurements, the estimates of the two expressions have their own errors and covariance; Chapter~\ref{ch:uncertainty} gives the theoretical propagation rule. Detailed measurement designs belong to the later books. The next theoretical step is the narrower thermodynamic connection to medium entropy under the accepted P4 assumptions.
